% Auto-generated from report/6.md — do not edit by hand.
\section{بحث مکانیکی و تبیین چرایی نتایج (\lr{Mechanistic Discussion})}\label{sec:6-1}


یافته‌های تجربی ارائه‌شده در فصل پنجم، پنج فرضیه علمی مطرح‌شده در فصل اول را تا حد زیادی تایید می‌کنند. واکاوی رفتار مدل‌ها در چهار بنچمارک زیستی، سه بینش بنیادین تئوریک پیرامون دینامیک یادگیری روی گراف‌های زیست‌مولکولی را آشکار می‌سازد.

\subsection{نقش عملگر پیوسته تخصیص منابع در مهار سوگیری هاب‌ها}\label{sec:6-1-1}


یکی از مهم‌ترین یافته‌های این پژوهش، تاثیر جایگزینی عملگر باینری $\text{\lr{sign}}(AA^T)$ با عملگر پیوسته تخصیص منابع ($W_2 = AD^{-1}A^T$) است. در گراف‌های زیستی که دارای توزیع درجه بی‌مقیاس هستند، گره‌های هاب (مانند داروهای پرعارضه یا پروتئین‌های چندمنظوره) به‌طور مصنوعی تعداد زیادی مسیر ۲-گام ایجاد می‌کنند.

عملگر باینری \lr{SkipGNN}، تمام این مسیرها را با وزن یکسان ($1.0$) لحاظ می‌کند که منجر به پدیده «اشباع سیگنال» در گره‌های هاب می‌شود. در مقابل، عملگر تخصیص منابع با تقسیم وزن هر مسیر بر درجه گره واسط ($1/d_k$)، به‌عنوان یک فیلتر پایین‌گذر توپولوژیک عمل کرده و سیگنال‌های اختصاصی (از طریق گره‌های واسط با درجه پایین) را تقویت می‌کند. این مکانیزم دلیل اصلی جهش عملکرد مدل \lr{AMS} در دیتاست متراکم \lr{DDI} (از $0.724$ به $0.912$ در آزمون سخت، معادل بهبود $+0.188$ واحد درصد) است، جایی که تداخلات دارویی از طریق مسیرهای چندگانه و پیچیده شکل می‌گیرند.

\subsection{حل بن‌بست هندسی دوقطبی با مسیرهای ۳-گام}\label{sec:6-1-2}


همان‌طور که در بخش ۲-۴-۱ اثبات شد، در گراف‌های دوقطبی (\lr{DTI} و \lr{GDI})، ماتریس $A^2$ ساختاری بلوکی-قطری دارد و فاقد هرگونه یال میان موجودیت‌های ناهم‌نوع است. مدل مرجع \lr{SkipGNN} با تکیه بر این ماتریس، در عمل دچار افت عملکرد تدریجی می‌شد (مشاهده‌شده در منحنی‌های یادگیری \lr{GDI} در شکل ۹).

مدل‌های \pcode{ThreeHopSkipGNN} و \pcode{AMS-SkipGNN} با تزریق صریح مسیرهای ۳-گام ($W_3 = W_2 D^{-1} A$)، این بن‌بست هندسی را شکستند. مسیر $d \to p' \to d' \to p$ (دارو \faarrow{} پروتئین واسط \faarrow{} داروی واسط \faarrow{} پروتئین هدف) اولین مسیر توپولوژیک غیرمستقیم برای دسترسی به تارگت کاندید است. موفقیت این مدل‌ها در دیتاست \lr{DTI} (کسب \lr{Hard AUPRC} برابر با $0.803$ برای \lr{3-hop} و $0.798$ برای \lr{AMS}) نشان می‌دهد که دسترسی انکودر به اطلاعات نوع مخالف، پیش‌شرط یادگیری در گراف‌های دوقطبی است.

\subsection{هم‌افزایی گیتینگ تطبیقی و دیکودر تانسوری تعاملی}\label{sec:6-1-3}


تحلیل نردبان ابلیشن (شکل ۲) یک الگوی قابل‌توجه را نشان داد: افزودن پرش وزن‌دار (گام ۱) و گیتینگ تطبیقی (گام ۲) به‌تنهایی و بدون اصلاح دیکودر، نه‌تنها بهبودی ایجاد نکرد، بلکه منجر به افت جزئی عملکرد (از $0.703$ به $0.686$ در \lr{DTI}) شد.

\textbf{تبیین مکانیکی:} پرش پیوسته \lr{RA} و گیتینگ، ظرفیت بازنمایی (\lr{Representation Capacity}) شبکه را به‌شدت افزایش می‌دهند و سیگنال‌های توپولوژیک غنی‌تری تولید می‌کنند. با این حال، دیکودر خطی الحاقی ($[E_u \parallel E_v]$) فاقد ظرفیت لازم برای پردازش این سیگنال‌های پیچیده و غیرخطی است. در نتیجه، شبکه دچار بیش‌برازش روی نویزهای فرکانس‌بالای گراف پرش می‌شود.

جهش اصلی در گام ۳ (افزودن دیکودر تانسوری ۴‌گانه) رخ داد. مولفه‌های ضربی ($E_u \odot E_v$) و فاصله‌ای ($|E_u - E_v|$) در این دیکودر، تقارن ذاتی تعاملات بیوشیمیایی (مانند مدل قفل‌وکلید) را به‌صورت صریح مدل‌سازی می‌کنند. این یافته نشان می‌دهد که در معماری‌های \lr{GNN}، ارتقای انکودر باید همواره با ارتقای متناسب ظرفیت دیکودر همراه باشد تا از پدیده «تنگنای اطلاعاتی» (\lr{Information Bottleneck}) جلوگیری شود.

\subsection{تمایز عملکرد بین دیتاست‌ها و کارآمدی \lr{SkipGATv2}}\label{sec:6-1-4}


یکی از یافته‌های مهم این پژوهش، تفاوت الگوی برتری مدل‌ها در دیتاست‌های مختلف است:

\begin{itemize}
\item \textbf{\lr{DDI} (گراف همگن متراکم):} مدل \lr{AMS} با $0.912$ برترین عملکرد را دارد، و برتری آن نسبت به سایر مدل‌های پیشنهادی کاملاً آشکار است.
\item \textbf{\lr{DTI} (گراف دوقطبی تنک):} سه مدل \lr{AMS}، \lr{SkipGATv2} و \lr{3-hop} با میانگین‌های $0.798$، $0.804$ و $0.803$ از نظر آماری هم‌سطح هستند (بازه‌های انحراف معیار هم‌پوشان).
\item \textbf{\lr{PPI} (گراف همگن مدولار):} \lr{SkipGATv2} با $0.778$ بهترین عملکرد را دارد که بهبود قابل‌توجه $+0.104$ واحد درصدی نسبت به \lr{AMS} نشان می‌دهد.
\item \textbf{\lr{GDI} (گراف دوقطبی کلان):} هوریستیک تحلیلی با $0.842$ بالاتر از تمام مدل‌های عصبی است.
\end{itemize}

\textbf{تحلیل \lr{SkipGATv2} در \lr{PPI}:} نکته مهم این است که برتری \lr{SkipGATv2} در \lr{PPI} \textbf{علی‌رغم} گام‌های بهینه‌سازی کمتر رخ داده است. این مدل در \lr{PPI} تنها ۳۱ گام \lr{Adam} در هر اپاک دریافت کرده (در مقایسه با ۲۵۵ گام \lr{AMS}، به دلیل بچ‌سایز ۱۰۲۴ در برابر ۱۲۸). این برتری با وجود گام‌های بهینه‌سازی کمتر، به‌عنوان شواهدی بر کارآمدی ذاتی مکانیزم توجه پویا در گراف‌های همگن مدولار تفسیر می‌شود، نه صرفاً نتیجه بودجه محاسباتی بیشتر.

\section{تحلیل پاشنه‌های آشیل و محدودیت‌های روش‌شناختی}\label{sec:6-2}


یک گزارش علمی معتبر موظف است محدودیت‌ها و شرایط مرزی مدل‌های خود را به‌صورت شفاف و بدون اغراق تبیین کند. تحلیل نتایج، چهار چالش و محدودیت مهم را آشکار می‌سازد که مسیرهای اصلی برای کارهای آینده هستند.

\subsection{چالش گراف‌های کلان و تنک: برتری هوریستیک در \lr{GDI}}\label{sec:6-2-1}


یکی از مهم‌ترین و صادقانه‌ترین یافته‌های این پژوهش، عملکرد خیره‌کننده هوریستیک تحلیلی \lr{Resource Allocation} در دیتاست ژن-بیماری (\lr{GDI}) است. در این گراف دوقطبی کلان ($19{,}783$ نود)، هوریستیک ۳-گام بدون هیچ پارامتر یادگیرنده‌ای، \lr{Hard AUPRC} برابر با $0.842$ کسب کرد که بالاتر از تمام مدل‌های عصبی (از جمله \lr{SkipGATv2} با $0.837$ و \lr{AMS} با $0.829$) است.

\textbf{تحلیل ریشه‌ای:} این پدیده که در ادبیات \lr{GNN}ها به عنوان «محدودیت سوگیری القایی» (\lr{Inductive Bias Limitation}) شناخته می‌شود، نشان می‌دهد که در گراف‌های به‌شدت تنک و کلان با ویژگی‌های هویتی، مکانیزم \lr{Message-Passing} در \lr{GNN}ها (حتی با وجود پرش‌های وزنی و گیتینگ) دچار پدیده «رقیق‌شدن سیگنال» (\lr{Signal Dilution}) می‌شود. هوریستیک \lr{RA} ذاتاً یک شمارشگر دقیق و بهینه مسیرهای ۳-گام است و در غیاب ویژگی‌های غنیِ گره‌ای (\lr{Node Features})، توپولوژی خالص را بهتر از یک شبکه عصبی ۲-لایه استخراج می‌کند.

\textbf{تبیین صادقانه:} این یافته نشان می‌دهد که معماری‌های عصبی برای غلبه بر هوریستیک‌ها در گراف‌های کلان، نیازمند تلفیق با ویژگی‌های معنادار بیوشیمیایی هستند تا صرفاً توپولوژی. به بیان دیگر، در این پژوهش با انتخاب آگاهانه ویژگی‌های هویتی (بخش ۱-۵-۲)، ما به‌طور عمدی مزیت ویژگی‌های بیوشیمیایی را حذف کردیم تا اثر خالص معماری را بسنجیم. این محدودیت به‌صراحت به‌عنوان مسیر کار آینده ذکر شده است.

\subsection{محدودیت ویژگی‌های هویتی (\lr{One-Hot}) و افت مکانیزم توجه در گراف کلان}\label{sec:6-2-2}


در دیتاست \lr{PPI}، مدل \lr{SkipGATv2} با کسب \lr{Hard AUPRC} برابر با $0.778$ بهترین عملکرد را در میان تمام مدل‌ها داشت که نشان‌دهنده کارآمدی مکانیزم توجه پویا در گراف‌های همگن با ساختار مدولار است. با این حال، همین مدل در دیتاست کلان \lr{GDI} دچار افت عملکرد شد و نتوانست از هوریستیک پیشی بگیرد.

\textbf{تبیین تئوریک:} این پدیده به «بیش‌تطبیق پارامتری توجه» (\lr{Attention Over-parameterization}) معروف است. زمانی که ویژگی‌های ورودی گره‌ها صرفاً بردارهای هویتی (\lr{One-Hot}) باشند، مکانیزم توجه در گراف‌های بسیار بزرگ فاقد سیگنال‌های معنادار اولیه برای محاسبه ضرایب توجه پایدار است. در نتیجه، مدل به‌جای یادگیری الگوهای ساختاری، دچار نویز محاسباتی می‌شود. این محدودیت مستقیماً ناشی از انتخاب روش‌شناختی این پژوهش مبنی بر ایزوله‌سازی اثرات توپولوژی (بخش ۱-۵-۲) است و در صورت افزودن ویژگی‌های پیوسته (مانند امبدینگ‌های توالی پروتئین)، برطرف خواهد شد.

\subsection{محدودیت‌های آماری و تعداد سیدها}\label{sec:6-2-3}


\textbf{تعداد سیدها:} به دلیل هزینه محاسباتی بالای پروتکل ارزیابی دوگانه (ساخت بانک منفی سخت، آموزش ۷ مدل، و ۴ دیتاست بر روی \lr{GPU T4})، آزمایش‌ها با ۳ سید تصادفی (۴۲، ۱۲۳، ۷) انجام شدند. این تعداد از نظر ادبیات \lr{GNN} (که معمولاً ۵ تا ۱۰ سید مرسوم است) کمتر از حد مطلوب است.

\textbf{آزمون‌های آماری:} در این گزارش، معیار تمایز آماری، عدم هم‌پوشانی بازه‌های میانگین $\pm$ انحراف معیار بوده است. آزمون‌های معناداری رسمی (مانند \lr{paired t-test} یا \lr{Wilcoxon signed-rank test}) به دلیل تعداد محدود سیدها انجام نشده‌اند. این محدودیت به‌ویژه در دیتاست \lr{DTI} اهمیت دارد، جایی که سه مدل پیشنهادی بازه‌های هم‌پوشان دارند و برتری یکی بر دیگری از نظر آماری قابل اثبات نیست.

\textbf{راهکار پیشنهادی:} در کارهای آینده، اجرای آزمایش‌ها با حداقل ۵ سید و گزارش \lr{p-value} رسمی برای هر مقایسه، ضروری است.

\subsection{محدودیت مقیاس‌پذیری}\label{sec:6-2-4}


بزرگ‌ترین گراف بررسی‌شده در این پژوهش حدود ۲۰ هزار نود داشت. تعمیم این معماری‌ها به گراف‌های میلیون‌نودی (مانند کل اینتراکتوم انسان) نیازمند بهینه‌سازی‌های سطوح پایین‌تر سخت‌افزاری (مانند \lr{FlashAttention} برای گراف‌ها) و تکنیک‌های نمونه‌برداری زیرگراف (\lr{Subgraph Sampling}) است که در این پروژه بررسی نشده‌اند.

الگوریتم \lr{Chunking} پیاده‌سازی‌شده برای \lr{SkipGATv2} (بخش ۴-۳-۳)، قله مصرف حافظه را از $O(|\mathcal{E}|)$ به $O(\text{\lr{chunk}})$ کاهش می‌دهد، اما زمان آموزش همچنان با تعداد یال‌های پرش به‌صورت خطی رشد می‌کند.

\section{نتیجه‌گیری و پاسخ به پرسش‌های پژوهش}\label{sec:6-3}


این پروژه با هدف بازبینی نقادانه و بازطراحی ساختاری مدل \lr{SkipGNN} در پیش‌بینی تعاملات زیست‌مولکولی اجرا گردید. نتایج به‌دست‌آمده به پنج پرسش بنیادین پژوهش به شرح زیر پاسخ می‌دهند:


\begin{table}[htbp]
\centering
\footnotesize
\caption{خلاصه پاسخ‌های پژوهش به پرسش‌های پنج‌گانه مطرح‌شده در فصل اول}
\label{tab:6-1}
\begin{tabularx}{\textwidth}{>{\hspace{0pt}}p{0.22\textwidth}>{\hspace{0pt}}X>{\hspace{0pt}}X}
\toprule
پرسش & فرضیه مطرح‌شده & نتیجه تجربی و وضعیت اثبات \\
\midrule
\textbf{\lr{RQ1}} (چگالی توپولوژیک پرش) & جایگزینی $\text{\lr{sign}}$ با پرش وزن‌دار \lr{RA} مانع از افت در منفی‌های سخت می‌شود. & \textbf{تایید مشروط:} پرش \lr{RA} به‌تنهایی کافی نیست، اما در ترکیب با دیکودر تانسوری (مدل \lr{AMS})، منجر به بهبود $+0.188$ واحد درصدی در \lr{DDI} و $+0.103$ واحد درصدی در \lr{DTI} نسبت به \lr{SkipGNN} شد. \\
\textbf{\lr{RQ2}} (دوگانگی ساختار دوقطبی) & مسیرهای ۳-گام بن‌بست گام زوج در گراف‌های دوقطبی را جبران می‌کنند. & \textbf{تایید:} مدل \lr{3-hop} با صراحت این بن‌بست را هدف قرار داد و در \lr{DTI} به \lr{Hard AUPRC} برابر با $0.803$ دست یافت (هم‌سطح با \lr{AMS} و \lr{SkipGATv2}). \\
\textbf{\lr{RQ3}} (ظرفیت دیکودر) & دیکودر تانسوری ۴‌گانه تعاملات غیرخطی متقارن را تفکیک می‌کند. & \textbf{تایید:} نردبان ابلیشن نشان داد که جهش اصلی عملکرد (از $0.686$ به $0.788$ در \lr{DTI}) دقیقاً در مرحله افزودن دیکودر تانسوری رخ می‌دهد. \\
\textbf{\lr{RQ4}} (پایداری فضای بازنمایی) & زیان \lr{InfoNCE} مانع از انحراف و چسبندگی امبدینگ‌های هاب می‌شود. & \textbf{تایید کیفی:} تحلیل بصری \lr{t-SNE} نشان داد که مدل \lr{Contrastive} دارای توزیع یکنواخت‌تر (\lr{Uniformity}) در فضای نهفته است، هرچند در معیارهای کمّی رتبه‌بندی، برتری مطلق نسبت به \lr{AMS} نداشت. \\
\textbf{\lr{RQ5}} (اعتبار پروتکل ارزیابی) & ارزیابی یکنواخت گمراه‌کننده است و ارزیابی سخت استاندارد واقعی است. & \textbf{تایید:} عملکرد \lr{SkipGNN} که در آزمون یکنواخت رقابتی به‌نظر می‌رسید ($0.929$ در \lr{DTI})، در آزمون سخت به‌شدت افت کرد ($0.695$)، که سوگیری پروتکل‌های پیشین را تایید می‌کند. \\
\bottomrule
\end{tabularx}
\end{table}


\textbf{پیام علمی محوری:} در پیش‌بینی تعاملات زیست‌مولکولی روی گراف، طراحی توپولوژیک عملگرهای پرش بر پایه نوع گراف (همگن در برابر دوقطبی) و استفاده از دیکودرهای تعاملی تانسوری، نقشی تعیین‌کننده ایفا می‌کنند. با این حال، در گراف‌های کلان و به‌شدت تنک، هوریستیک‌های تحلیلی همچنان رقبای جدی برای مدل‌های عصبیِ مبتنی بر توپولوژی خالص هستند.

\textbf{شفاف‌سازی مهم:} مدل \lr{AMS} عمیق‌ترین اعتبارسنجی ساختاری (نردبان ابلیشن کامل) را در میان چهار مدل دارد و در دیتاست \lr{DDI} به‌وضوح برترین مدل عصبی است. با این حال، از نظر امتیاز خام \lr{Hard AUPRC}، در دیتاست‌های \lr{DTI}، \lr{PPI} و \lr{GDI} با \lr{SkipGATv2} هم‌سطح یا در محدوده انحراف معیار آن قرار دارد. این تنوع عملکردی، غنای یافته‌های پژوهش است و نشان می‌دهد که یک معماری واحد نمی‌تواند در تمام ساختارهای گرافی برتری مطلق داشته باشد.

\section{افق‌های پژوهشی و کارهای آینده}\label{sec:6-4}


بر پایه یافته‌ها و محدودیت‌های شناسایی‌شده در این پروژه، مسیرهای پژوهشی زیر به‌عنوان توسعه‌های آتی پیشنهاد می‌گردد:

\subsection{تزریق ویژگی‌های چندوجهی بیوشیمیایی پیش‌آموزش‌دیده}\label{sec:6-4-1}


مهم‌ترین محدودیت این پژوهش، اتکا به ویژگی‌های هویتی (\lr{One-Hot}) بود. در کارهای آینده، افزودن ویژگی‌های پیوسته استخراج‌شده از مدل‌های زبانی بزرگ زیستی می‌تواند پاشنه آشیل مدل‌ها در گراف‌های کلان (مانند \lr{GDI}) را برطرف کند. این ویژگی‌ها شامل:

\begin{itemize}
\item \textbf{برای مولکول‌های دارویی:} اثرانگشت‌های ساختاری مورگان \cite{landrum2016rdkit} و امبدینگ‌های مدل \lr{ChemBERTa} \cite{chithrananda2020chemberta} برای کپچر کردن خصوصیات شیمیایی.
\item \textbf{برای توالی‌های پروتئینی:} امبدینگ‌های لایه‌های آخر مدل‌های زبانی پروتئین نظیر \lr{ESM-2} \cite{lin2023esm2} یا \lr{ProtTrans} \cite{elnaggar2022prottrans}.
\item \textbf{برای ژن‌ها:} امبدینگ‌های \lr{Gene Ontology} یا الگوهای بیان ژن.
\end{itemize}

این تلفیق می‌تواند افت \lr{SkipGATv2} در گراف‌های کلان (بخش ۶-۲-۲) را جبران کند.

\subsection{معماری ترکیبی ۳-گام و یادگیری خودنظارتی}\label{sec:6-4-2}


ترکیب عملگرهای مسیر طول فرد در مدل \pcode{ThreeHopSkipGNN} با قید یکنواختی کروی در \pcode{ContrastiveSkipGNN} به‌منظور ایجاد مدلی که همزمان فیزیک گراف دوقطبی را بشناسد و فضا را در برابر بیش‌برازش هاب‌ها تنظیم نماید. این ترکیب می‌تواند به عنوان یک معماری هیبریدی در گراف‌های دوقطبی تنک عملکرد بهینه ارائه دهد.

\subsection{تعمیم به گراف‌های دانش ناهمگن چندنوعی}\label{sec:6-4-3}


گسترش عملگرهای پرش چندمقیاسی به شبکه‌های دارای ده‌ها نوع موجودیت همزمان (دارو، ژن، بیماری، مسیر بیولوژیکی، فنوتیپ و عوارض جانبی) در قالب شبکه‌های عصبی گراف رابطه‌ای (\lr{Relational GNNs}) \cite{schlichtkrull2018rgcn} و پیش‌بینی چندتکلیفی (\lr{Multi-task Learning}).

\subsection{توسعه الگوریتم‌های شتاب‌دهی سخت‌افزاری}\label{sec:6-4-4}


بازنویسی لایه‌های اتنشن تنک (مانند \pcode{SparseGATv2Layer}) با کرنل‌های ادغام‌شده \lr{Triton}/\lr{CUDA} (مشابه \lr{FlashAttention}) جهت حذف کامل سربار محاسباتی و امکان آموزش روی گراف‌های بالاتر از ۱۰۰ هزار گره بدون نیاز به تکنیک‌های \lr{Chunking}.

\subsection{اعتبارسنجی آماری با سیدهای بیشتر}\label{sec:6-4-5}


اجرای آزمایش‌های نهایی با حداقل ۵ تا ۱۰ سید تصادفی و گزارش آزمون‌های معناداری رسمی (\lr{paired t-test} یا \lr{Wilcoxon signed-rank}) برای هر مقایسه مدل. این کار به‌ویژه برای دیتاست‌هایی مانند \lr{DTI} که در آن بازه‌های انحراف معیار هم‌پوشان هستند، ضروری است.

\subsection{معیارهای کمّی برای فضای نهفته}\label{sec:6-4-6}


افزودن معیارهای کمّی مانند \lr{Silhouette Score}، \lr{Alignment} و \lr{Uniformity} \cite{wang2020uniformity}، یا دقت طبقه‌بندی \lr{k-NN} روی فضای نهفته به‌عنوان مکمل تحلیل بصری \lr{t-SNE}. این معیارها می‌توانند تبیین دقیق‌تری از کیفیت فضاهای نهفته مدل‌های مختلف ارائه دهند.

